\section{Related Work}
\label{sec:related-work}
\label{sec:related}

\subsection{Benchmarks for LLM Code Agents}

\textbf{Bug-fix benchmarks.}
SWE-bench~\cite{jimenez2024swebench} and its human-verified subset SWE-bench~Verified~\cite{chowdhury2024swebenchverified} established the modern standard for evaluating LLM coding agents on real GitHub bug-fix patches.
Subsequent work extends the formula to multi-file repair (SWE-Bench~Pro~\cite{swebenchpro2025}), polyglot corpora (SWE-Bench~Live~\cite{swebenchlive}, Multi-SWE-Bench~\cite{multiswebench}), and non-trivial test-driven patches (TDD-Bench-Verified~\cite{tddbench}).
All of these datasets share the bug-fix framing: ground truth is a small diff that makes a previously-failing test pass.
Our work targets the complementary regime of \emph{feature implementation}, where tasks include multi-file refactors, new subsystems, and cross-module protocol changes---none of which are reducible to a single failing test.

\textbf{Feature-level benchmarks.} The closest relatives to our study are FeatBench~\cite{featbench2025} and SWE-EVO~\cite{sweevo}, which introduce feature-addition evaluations, and CR-Bench~\cite{crbench} for code-review-grounded evaluation.
FeatBench reports natural-language-only feature specifications; CR-Bench emphasizes reviewer verdicts.
Our CapBench extension is closest in spirit but differs in two ways: (i)~we include a proprietary-partner subset (Huawei~CANN, MindSpeed) that represents accelerator-stack engineering absent from prior open benchmarks; and (ii) we measure \emph{implementation equivalence} against merged ground truth, which explicitly penalizes under-specified diffs.

\textbf{Industrial-context benchmarks.}  Passerine~\cite{mathai2024passerine} evaluates LLM-assisted program repair at Google; ProdCodeBench~\cite{prodcodebench2024} benchmarks production-grade code completion; TestGen-LLM at Meta~\cite{alshahwan2024testgen} studies unit-test generation at scale. Unlike these works, our benchmark targets \emph{feature implementation}---multi-file, behavior-changing modifications with a merged pull-request as ground truth---rather than repair or completion.

\subsection{LLM Coding Agents and Harness Design}

\textbf{Model families.} The agents we evaluate embed different frontier models (Claude~4.6 Opus, GPT-5.4, GLM-4.5, and Cursor's proprietary routing system). Prior comparative studies such as MultiAgent~\cite{multiagent2024} and SWE-bench agents~\cite{jimenez2024swebench} evaluate agents on bug-fixing traces; we extend this to feature implementation across industrial repositories.

\textbf{Agent architectures.} Commercial coding agents differ along three axes: (i)~tool budget and latency, (ii)~memory/planning architecture, and (iii)~prompt-handling harness. Prior harness-engineering work~\cite{harness2025,agentless,confucius} has shown that simple pipeline-level choices (file-level planning, edit localisation, self-reflection) dominate raw model capability on SWE-bench Verified. Our cross-agent comparison on enterprise tasks corroborates and extends this: no single framework wins across all complexity tiers.

\paragraph{Industrial deployment studies.}
LogSage~\cite{logsage}, Passerine~\cite{mundler2024passerine}, and Meta's AutoBug~\cite{autobug} report deployments at Huawei, Google, and Meta respectively, each on a narrower task class (CI/CD triage, test-bug fix, auto-remediation).
Our work complements these by evaluating the \emph{generic} feature-implementation task across four commercial agents and twelve industrial tasks plus fifty open-source controls.

\paragraph{Complementary evaluation perspectives.}
TestGen-LLM~\cite{testgenllm} and CodeVisor~\cite{codevisor} evaluate test-centric code generation; ProdCodeBench~\cite{prodcodebench2024} and RepoFusion~\cite{repofusion} evaluate repository completion. These targets are complementary to ours: successful feature implementation presupposes the test synthesis and code-completion capabilities they measure in isolation.
